\section{November 2026 expansion: evidence and reading guide}
\label{sec:expansion-guide}
This edition adds explanations requested in the supplied ``To do list for the
next 11 days''. The reference report is the unsuffixed
\texttt{23539195\_BrauneDR\_FINALREPORT.pdf} in Downloads, created by Prism on
3 November 2026, 108 PDF pages. Its printed page numbers differ from PDF page
indices. The original per-function chapters are retained. New sections are
inserted within the corresponding chapters, followed by a separate chapter on
measurement and verification. This is a study manual, not an amended report or
a record of new experiments.

\paragraph{What the evidence labels mean.}
\textbf{Derived} means a result follows from explicitly stated assumptions and
algebra; it does not establish that those assumptions describe the hardware.
\textbf{Published method} means a cited source supports the method, not every
project coefficient. \textbf{Reported observation} means the final report or
archived record contains it, without a new execution here.
\textbf{Project heuristic} means an engineering threshold or preference whose
exact optimum has not been established in this edition.
\textbf{Unresolved} means the available evidence is contradictory or insufficient.
Arithmetic examples in the added sections are explanatory examples unless
explicitly identified as reported observations. They are not new measurements.

\paragraph{Version boundary.}
This manual originated against CleanPipeline2; the report and subsequent
CleanPipeline3 implementation are not identical snapshots. Read-only inspection
of \path{Software/CleanPipeline3/segment.py}, \texttt{recognise.py},
\texttt{style.py}, \texttt{synthesise.py}, \texttt{reproduce.py},
\texttt{train.py} and \texttt{paths.py} informs the additions. A source comment
can explain intent but cannot prove a performance claim. Checkpoint metadata
and as-run logs outrank defaults for historical settings. None of these modules
was imported, executed or edited for this expansion. The separate agent's
ongoing coefficient sweeps are outside this edition's evidence boundary.

\paragraph{Where the requested subjects are explained.}
\begin{description}[style=nextline]
\item[Segmentation, Section~\ref{sec:expand-segment}.] Integral images;
illumination ratio; histogram/Otsu derivation; orientation and deskew;
run removal, pitch, 6\% peaks and $0.6P$ spacing; seam cost and dynamic
programming; closing radius and component ownership; contrast and darkness;
complete photographed-page-to-line algorithm and interpretation of seam plots.
\item[Recognition, Section~\ref{sec:expand-recognise}.] Input sizing;
convolutional map $Z$ and feature sequence; channel counts and receptive fields;
normalisation, ReLU, sigmoid, pooling, LSTM and BiLSTM; CTC, collapse,
greedy decoding, forced alignment and frame coordinates; writer head,
$K$ classes, dropout; stroke normalisation and its recorded sweep.
\item[Profiles, Section~\ref{sec:expand-style}.] Body-band guess;
45\%, percentile and cut-cost terms; thinning, dilation, graph tracing,
75-degree continuation, Douglas--Peucker and $0.6$ pixel tolerance;
coordinate shear; width gates; $12\times12$ descriptor, prior, consensus,
medoid versus prototype; feature measurement and connectivity; full extraction
algorithm and its failure modes.
\item[Synthesis, Section~\ref{sec:expand-synth}.] Coordinates, layout,
drift persistence and clearance; automatic anchors; $0.45$ connection gate,
$0.8$ x-height gap; joint character/writer score and harmonic mean;
candidate search and accepted repair; effects of removing each mechanism.
\item[Motion and electronics, Section~\ref{sec:expand-hardware}.] Torque,
units, interpolation assumptions, resolution and quantisation;
A4988 reference voltage; BC337 and resistors; every documented GPIO and driver
connection; pen state machine, debounce, calibration, G-code, Bresenham,
trajectory bounds, timing and $0.6$ mm joins.
\item[Training, Section~\ref{sec:expand-train}.] Seeded 15\% holdout,
20\% personal sampling, text/writer settings, loss scaling, backward
derivatives, optimisers, clipping, schedules, augmentation, complete training
logic and historical epoch/time evidence.
\item[Verification, Chapter~\ref{ch:expand-verification}.] First-principles
versus Torch comparisons; finite differences and recorded arithmetic;
Appendix A.0.1 vote; CER/Levenshtein, Wilson intervals, geometric registration,
point-set F-score, natural-variation ratio, direct lengths and timing;
capture/seam records, plot-reading guides and coefficient evidence register.
\end{description}

\paragraph{How sources should be read.}
New sections distinguish published definitions from original project-specific
reasoning. Mathematical derivations and small constructed examples in these
sections are supplied to make the reasoning checkable. A reference to Otsu,
CTC or Douglas--Peucker establishes an algorithm's origin; it cannot justify
the project's $0.55$, $0.6$ or $1.15$ constants. The bibliography supplies
persistent links for the additions, and Section~\ref{sec:source-audit} states
which original references could not be independently checked. This edition
does not certify every reference in the final report.
